\chapter{Control Flow}
\label{chap:control-flow}

\chapterguide
  {You should be able to construct logical expressions.}
  {use conditional statements, \texttt{switch()}, loops, \texttt{break}, and \texttt{next}.}

\section{Selection statements}

Conditional statements select which code block runs.

\subsection{Simple \texttt{if}}

\begin{lstlisting}[style=dsr]
x <- 30L

if (is.integer(x)) {
  print("X is an Integer")
}
\end{lstlisting}

The body runs only when the condition is TRUE.

\subsection{\texttt{if...else}}

\begin{lstlisting}[style=dsr]
x <- c("what", "is", "truth")

if ("Truth" %in% x) {
  print("Truth is found")
} else {
  print("Truth is not found")
}
\end{lstlisting}

Character comparison is case-sensitive: \texttt{"Truth"} and \texttt{"truth"} differ.

\subsection{\texttt{if...else if...else}}

\begin{lstlisting}[style=dsr]
x <- c("what", "is", "truth")

if ("Truth" %in% x) {
  print("Truth is found the first time")
} else if ("truth" %in% x) {
  print("truth is found the second time")
} else {
  print("No truth found")
}
\end{lstlisting}

Conditions are checked in order until one succeeds.

\section{\texttt{switch()}}

With a numeric selector, \texttt{switch()} returns the matching position:

\begin{lstlisting}[style=dsr]
x <- switch(
  3,
  "first",
  "second",
  "third",
  "fourth"
)
print(x)
\end{lstlisting}

Selector \texttt{3} returns \texttt{"third"}.

\section{Repetition statements}

R provides \texttt{repeat}, \texttt{while}, and \texttt{for} loops. They differ
only in where the decision to continue is made, which is easier to see drawn
than described.

\begin{utpfig}[The three loops]{Where each loop tests. \texttt{repeat} has no
test of its own and must be broken from inside. \texttt{while} tests before
every pass, so its body may never run. \texttt{for} runs once per element and
cannot loop forever.}{fig:loops}
\begin{tikzpicture}[x=1cm,y=1cm]
% ---- repeat
\node[tagb] at (1.3,1.35) {\texttt{repeat}};
\node[fcstart] (r0) at (1.3,0.75) {start};
\node[fcbox,minimum width=17mm] (r1) at (1.3,-0.15) {body};
\node[fcdiamond] (r2) at (1.3,-1.35) {break?};
\node[fcstart] (r3) at (1.3,-2.5) {end};
\draw[fcedge] (r0)--(r1);
\draw[fcedge] (r1)--(r2);
\draw[fcedge] (r2)--node[fclabel,right]{yes}(r3);
\draw[fcedge] (r2.west) -- ++(-1.0,0) |- node[fclabel,left,pos=0.25]{no} (r1.west);

% ---- while
\node[tagb] at (5.3,1.35) {\texttt{while}};
\node[fcstart] (w0) at (5.3,0.75) {start};
\node[fcdiamond] (w1) at (5.3,-0.25) {cond?};
\node[fcbox,minimum width=17mm] (w2) at (5.3,-1.45) {body};
\node[fcstart] (w3) at (7.3,-0.25) {end};
\draw[fcedge] (w0)--(w1);
\draw[fcedge] (w1)--node[fclabel,right]{true}(w2);
\draw[fcedge] (w1)--node[fclabel,above]{false}(w3);
\draw[fcedge] (w2.west) -- ++(-1.0,0) |- (w1.west);

% ---- for
\node[tagb] at (10.2,1.35) {\texttt{for}};
\node[fcstart] (f0) at (10.2,0.75) {start};
\node[fcdiamond] (f1) at (10.2,-0.25) {more?};
\node[fcbox,minimum width=19mm] (f2) at (10.2,-1.45) {body with\\ next element};
\node[fcstart] (f3) at (12.3,-0.25) {end};
\draw[fcedge] (f0)--(f1);
\draw[fcedge] (f1)--node[fclabel,right]{yes}(f2);
\draw[fcedge] (f1)--node[fclabel,above]{no}(f3);
\draw[fcedge] (f2.west) -- ++(-1.0,0) |- (f1.west);
\end{tikzpicture}
\end{utpfig}

\subsection{\texttt{repeat} loop}

\begin{lstlisting}[style=dsr]
v <- c("Hello", "loop")
cnt <- 2

repeat {
  print(v)
  cnt <- cnt + 1

  if (cnt > 5) {
    break
  }
}
\end{lstlisting}

\texttt{repeat} requires an internal exit such as \texttt{break}.

\subsection{\texttt{while} loop}

\begin{lstlisting}[style=dsr]
v <- c("Hello", "while loop")
cnt <- 2

while (cnt < 7) {
  print(v)
  cnt = cnt + 1
}
\end{lstlisting}

\texttt{while} checks its condition before each iteration; the update prevents an infinite loop.

\subsection{\texttt{for} loop}

\texttt{for} iterates over vectors or lists:

\begin{lstlisting}[style=dsr]
v <- LETTERS[1:4]
for (i in v) {
  print(i)
}

for (x in 1:10) {
  print(x)
}

fruits <- list("apple", "banana", "cherry")
for (x in fruits) {
  print(x)
}
\end{lstlisting}

\section{Controlling a loop with \texttt{break} and \texttt{next}}

The two keywords are often confused because both interrupt the body. They stop
different things.

\begin{utpfig}[break against next]{\texttt{break} leaves the loop entirely;
\texttt{next} abandons the current iteration and goes back for the following
one. Neither affects an enclosing loop.}{fig:break-next}
\begin{tikzpicture}[x=1cm,y=1cm]
% the loop body, as a shaded region
\begin{scope}[on background layer]
\fill[UTPPaleGray,rounded corners=2mm,draw=UTPMuted!50]
  (-0.05,-3.05) rectangle (4.6,0.45);
\end{scope}
\node[tagb,anchor=north west] at (0.1,0.35) {loop body};

\node[fcstart] (top) at (2.25,-0.25) {next iteration starts};
\node[fcbox,minimum width=24mm] (s1) at (2.25,-1.1) {statement};
\node[fcdiamond] (t) at (2.25,-1.95) {test};
\node[fcbox,minimum width=24mm] (s2) at (2.25,-2.75) {statement};

\draw[fcedge] (top)--(s1);
\draw[fcedge] (s1)--(t);
\draw[fcedge] (t)--node[fclabel,right,pos=0.45]{otherwise}(s2);

% next: back to the top of the body
\draw[-{Stealth[length=2mm]},semithick,draw=UTPTeal]
  (t.west) -- ++(-1.55,0)
  |- node[fclabel,left,pos=0.25,text=UTPTeal]{\texttt{next}} (top.west);

% break: out of the body altogether
\node[fcstart,draw=UTPRed,fill=UTPPaleRed] (out) at (6.9,-1.95) {after the loop};
\draw[-{Stealth[length=2mm]},semithick,draw=UTPRed]
  (t) -- node[fclabel,above,text=UTPRed]{\texttt{break}} (out);

\node[tag,anchor=west,text width=4.3cm] at (5.4,-0.7)
  {\textcolor{UTPTeal}{\texttt{next}} skips the rest of \emph{this} pass.};
\node[tag,anchor=west,text width=4.3cm] at (5.4,-2.9)
  {\textcolor{UTPRed}{\texttt{break}} skips every remaining pass.};
\end{tikzpicture}
\end{utpfig}

\texttt{break} exits the current loop:

\begin{lstlisting}[style=dsr]
fruits <- list("apple", "banana", "cherry")

for (x in fruits) {
  if (x == "cherry") {
    break
  }
  print(x)
}
\end{lstlisting}

\texttt{next} skips the rest of the current iteration:

\begin{lstlisting}[style=dsr]
for (x in fruits) {
  if (x == "banana") {
    next
  }
  print(x)
}
\end{lstlisting}

\section{Conditions inside loops}

The dice example combines iteration, a comparison, and an \texttt{if...else} block:

\begin{lstlisting}[style=dsr]
dice <- 1:6

for (x in dice) {
  if (x == 6) {
    print(paste("The dice number is", x, "Yahtzee!"))
  } else {
    print(paste("The dice number is", x, "Not Yahtzee"))
  }
}
\end{lstlisting}

This pattern iterates, tests, and produces a conditional result.

\section{Nested loops}

The final Lab 5a example places one \texttt{for} loop inside another:

\begin{lstlisting}[style=dsr]
adj <- list("red", "big", "tasty")
fruits <- list("apple", "banana", "cherry")

for (x in adj) {
  for (y in fruits) {
    print(paste(x, y))
  }
}
\end{lstlisting}

The nested loops generate every adjective--fruit pair.

\begin{quickcheck}
\begin{enumerate}
  \item What must be TRUE for the body of a simple \texttt{if} statement to run?
  \item Why does the first branch of the handout's \texttt{"Truth"} example fail while the second succeeds?
  \item How does a \texttt{repeat} loop terminate in the supplied pattern?
  \item What is the difference between \texttt{break} and \texttt{next} in the fruit examples?
  \item How many adjective-fruit combinations are generated by three adjectives and three fruits?
\end{enumerate}
\end{quickcheck}

\begin{chapterrecap}
\begin{itemize}
  \item Conditionals and \texttt{switch()} select code paths.
  \item \texttt{repeat}, \texttt{while}, and \texttt{for} express different repetition patterns.
  \item \texttt{break} exits a loop; \texttt{next} skips one iteration.
\end{itemize}
\end{chapterrecap}
